Measurement design is an argument of the long-term-effect efficiency problem, not a harmless preprocessing choice. A full-catalogue pilot ranks the finite admissible designs by their profiled bridge-gradient value and learns the square-root source allocation. The allocation criterion has the exact curvature identity
\[
 F_d(\widehat\alpha)-\mathcal V_d^*(P)
 =\frac{\mathcal V_d^*(P)(\widehat\alpha-\alpha_d^*(P))^2}
 {\widehat\alpha(1-\widehat\alpha)},
\]
and integer purchasing adds only an explicit nonnegative reciprocal-floor remainder, so allocation estimation is second order. Prediction screening has no comparable guarantee. At the binary center \(P_{\chi,\epsilon}\), for every finite \(\chi>1\), it uniquely chooses a design that is costlier in each source by \(5(\chi-1)/2\), while both efficiency ratios exceed \((\epsilon^{-2}+14)/726\). Uniformly over \(\chi>1\), the sharp first-order elbow is
\[
 \lim_{\epsilon\downarrow0}\epsilon^2
 \inf_{\chi>1}\rho_{\mathrm{pred}}(P_{\chi,\epsilon})
 =\frac1{726},
\]
with the same limit under the IKM normalization. The arm-normalized value concerns the identified bridge contrast on the full class; absolute causal-efficiency language requires common identification and the fixed-system IKM premises.